% !TEX program = xelatex
\documentclass{article}
\usepackage{geometry}\geometry{margin=23mm}
\usepackage{fontspec}\setmainfont{Times New Roman}
\usepackage[french]{babel}
\begin{document}
\title{La planche de mesure}\author{}\date{14 mai 2026}\maketitle
\begin{abstract}Deux repères indiquent la position initiale de la planche.\end{abstract}
\tableofcontents
\section{Les repères}\begin{quote}« Lire d’abord le repère inférieur. »\end{quote}
\begin{figure}[ht]\centering\fbox{\rule{0pt}{8mm}\rule{33mm}{0pt}}\caption{Position de la planche}\end{figure}
\section{Hauteur}
\begin{table}[ht]\centering\begin{tabular}{cr}Repère&Hauteur (mm)\\\hline 1&12\\2&18\end{tabular}\caption{Hauteur des repères}\end{table}
\begin{thebibliography}{1}\bibitem{record}\emph{Carnet de mesure}, 2025, \chaptername~2.\end{thebibliography}
\end{document}
